\documentclass[12pt]{article}

\usepackage{a4wide}
\usepackage[latin1]{inputenc}
\usepackage{tikz-er2}


\begin{document}

\title{Movinator 1.0\\Specification}
\author{P�vel Calado}
\maketitle

{\Huge\textbf{CANCELADO!!!}}\\
J� existem ``movie social networks'', que, no fim de contas, � o que eu �a
fazer. Porra!

\section{Requirements}
\label{sec:requirements}

\subsection{Data Requirements}
\label{sec:requirements-description}

\begin{itemize}
\item The database stores movies. Movies have an id, year, an original title,
    and can have several titles. Each title has a content and a language. We
    need to know which user inserted wich movie in the database.
\item A movie can have one or more directors. A director has an id and a name.
\item The database also stores users. These have a login, a name, an id, a
    password, and an email address.
\item Users watch movies, on a given date. To any watched movie, they assign a
  score (0---awfull, through 5---masterpiece).
\item Users also assign tags to movies. A tag has an id and a description.
\item The database also stores critics. A critic has an id and a name.
\item Critics review movies, on a given date. To any reviewed movie they assign
    a score (0---awfull, through 5---masterpiece). The review can have an URL.
\item When critics review movies, they do it for a periodical. A periodical has
    a name, an id, and a URL.
\end{itemize}

\subsection{Functional Requirements}
\label{sec:funct-requ}

\begin{itemize}
\item All ids are numbers generated by the system.
\item Users can login to the system using their email and password.
\item Users insert movies, with the corresponding information, and can create
    and assign tags to movies.
\item Users can also insert critics and their reviews and periodicals.
\item Users can insert, delete and edit any information they also inserted.
\item When something is deleted, if no one else is using it, it is deleted from
    the system.
\item When they are inserting a movie, the system will look for movies with a
    similar title and suggest them, to avoid having the same movie repeated in
    the database.  When suggesting a movie, the system shows its year,
    directors, and how many users have reviewd that movie.
\item The same happens when the user is inserting a critic, a periodical, a
    tag, or a director.
\item For critics, the name, periodicals and number of movies is shown. For
    periodicals, the name, URL, and number of movies is shown. For tags, the
    number of the number of movies is shown. For directors, the number of
    movies directed is shown. The suggestions are sorted by number of movies.
\item When assigning reviews to movies, the user sees a list of the critics she
    inserted. She can also type the name of a critic, or use a suggestion, as
    above.
\item The system can reccomend movies.
\item The system should be available through the web.
\end{itemize}


\subsection{E-R Schema}
\label{sec:e-r-schema}


\begin{center}
    \begin{tikzpicture} [node distance=7em]
        \draw (0,0) grid (10,10);

    
    \end{tikzpicture}
\end{center}

\end{document}
